\section{产物的结构与组成表征}
\label{sec:validation-endpoints}

图~\ref{fig:route-matrix}\CJKecglue{}右侧列出的四类表征分别针对不同问题，不能相互替代：单晶或先进衍射方法确认结构模型，PXRD\slash Rietveld判断块体相组成，化学分析和显微观察检查实际元素比与污染，高温实验界定特定气氛和热历史下的稳定范围。表~\ref{tab:synthesis-conditions-b}\CJKecglue{}的“产物与表征”栏据此限定每条实验记录可支持的结论。缺少某类数据时，仅对相应结论保持谨慎，不能以其他方法补足。

\subsection{结构模型}

单晶X射线衍射可将目标相的判断落实到实测晶体的结构模型；目前的直接研究正是据此鉴定$\mathrm{Ba_5Y_{12}Zn[O(SiO_4)]_8}$\CJKecglue{}\cite{ababaikeri2024ba5y12zn}。在无Zn的相关Ba--Y硅酸盐中，单晶精修揭示了非化学计量调制结构，cRED、同步辐射PXRD与中子精修的组合确认了独立的四方Ba--Y相区\cite{yamane2024synthesis,gulay2024navigation}。其他单晶或同步辐射研究确立了各自Ba--Y、Ba--Zn--Si化合物的结构，但结论仅适用于相应化合物\cite{motozawa2022bay16si4o33,zou2021crystal}。单晶结构能证明被测晶体与精修模型一致，却不能单独证明整批样品相纯，也不能将名义投料组成等同于实测位点占据；局域谱学至多补充多晶型辨析，不提供目标相的直接相形成证据\cite{becerro2004revisiting}。

\subsection{块体相组成}

PXRD\slash Rietveld面向块体样品，重点分析相分数、未指标化衍射峰以及相同热历史下的晶型差异。Ba--Y陶瓷研究通过粉末衍射发现了第二相，独立相区研究将同步辐射粉末数据与其他衍射方法结合以确认相组成\cite{yamane2024microstructure,gulay2024navigation}；Ba--Zn--Si研究也利用X射线、中子或同步辐射数据区分块体晶型与结构模型\cite{lin1999phase,zou2021crystal}。因此，PXRD\slash Rietveld可支持块体相组成和杂相判断，但峰位相似或主相出现均不能自动推出相纯度，也不能替代单晶位点占据判断。对仅报道固相转变或热致变化的研究，本文仅借鉴其表征方法，不将题名中的“优化”视为目标相判据\cite{kerstan2012thermal,cai2024optimized}。

\subsection{实际组成与污染}

名义配比、局部测点组成和样品整体组成需要分别记录。EDS、EPMA或ICP可检查晶粒间的成分差异、外来相和处理介质带入。无Zn的Ba--Y陶瓷研究将第二相与玛瑙来源的$\mathrm{SiO_2}$污染同时纳入分析\cite{yamane2024microstructure}；独立Ba--Y相区研究以EDS配合多种衍射方法约束组成和结构\cite{gulay2024navigation}；Ba\slash Sr--Zn--Si相关化合物则以单晶结构和EDS联合确认自身的组成关系\cite{thieme2015ba1}。这些方法仅能说明测点或取样体积内检测到的元素，不能凭少量测点判断整批样品的均匀性，也不能将烧前名义配比直接改写为烧后化学式。显微、热学和光谱联合测试可为掺杂或晶体生长提供过程参照，但不承担目标相的组成结论\cite{tejas2024structural,pang2005study}。

\subsection{高温稳定性}

HT\nobreakdash-XRD、热分析、热膨胀测量和多气氛实验考察材料在特定热历史或气氛下的变化，不能替代室温结构鉴定。BaZn$_2$Si$_2$O$_7$的低温相与高温相由中子与X射线衍射区分，说明原位高温相不能等同于冷却后产物\cite{lin1999phase}；Zn\slash Co及Ba\slash Sr--Zn--Si\slash Ge相关研究进一步将取代组成与热稳定性或相边界联系起来\cite{kerstan2013bazn2si2o7,thieme2016negative}。相关陶瓷在特定还原气氛中的耐受性仅可作为气氛参照，不能证明目标相具有相同的稳定性\cite{zou2019anti}。高温实验可界定已发表体系的转变、分解或气氛耐受范围，但原位信号不能单独证明冷却后样品相纯；缺少相应数据时，稳定性结论保持未确定。

四类表征共同限定结论的适用范围：单晶或先进衍射确认结构模型，粉末衍射判断块体相组成，化学分析检验实际元素比与污染，高温实验界定热历史与气氛边界。只有将这些结果与同一批次的合成条件对应，才能讨论目标相是否已复现，以及相关工艺可外推的范围。
